% Discussion Section v3 — 2026-06-08 (Round 1 revised after adversarial + writing review)
% Reframe: CVD color deficit = individual-specific GEOMETRIC DISTORTION of
%   cortical color representation. RDM = structure, LOCO = functional consequence.
%   Neural + behavioral measurements JOINTLY ground an individualized correction.
% Spine: 3 contributions — C1 structural/geometric distortion (RQ1, ¶2);
%   C2 individualized filter from cortical structural info (¶3-5);
%   C3 performance superiority (Phase 3 forward-looking, ¶6).
% Removed from v2: upstream-input-rejection paragraph; detection-correction
%   "divergence" falsifier + §S16 cosine statistic.
% Brief / evidence pack / constraints: discussion_structure_v3.md

\section{Discussion}
\label{sec:discussion}

Two CVD participants, one deutan and one protan, showed not a uniform loss of signal but a structured geometric distortion of the cortical color representation. Both selectively lost continuous-hue interpolation and showed elevated representational disparity from healthy controls, while still discriminating all eight color hues at every early-visual ROI. We modeled it as a hue rotation about two color-space axes, the subtype confusion axis and the S-cone axis. The former is the classical direction of CVD color confusion, the latter where the interpolation deficit was concentrated. We fit this model independently per participant, from behavioral discrimination thresholds and cortical representational geometry. We then recovered an exact stimulus-space pre-image for all eight hues in both cases, the per-person correction filter built from each individual's own cortical color geometry. In a first two-person evaluation, this filter's effect on cortical hue interpolation was inconsistent: it moved the deutan participant partway toward the healthy-control reference but did not restore the protan participant. Behaviorally, the individualized filter reduced each participant's discrimination errors and, unlike the deployed accessibility filter, introduced no new ones. Whether it outperforms that filter could not be confirmed.

\subsection{A geometric distortion of cortical color representation}
In both CVD participants, the deficit appeared as a structured distortion of cortical color geometry rather than a uniform attenuation of the chromatic signal. Representational geometry treats the pairwise distance structure among stimuli as an encoded relationship in each cortical region \cite{kriegeskorte2008, kriegeskorte2019}. Two complementary measurements expose this distortion, each at the stage where its signal exists. The representational dissimilarity matrix (RDM) shows \emph{which} pairwise distances deviate from the healthy-control (HC) geometry: elevated disparity appeared at different areas in the two participants --- V2 in the deutan and V1 in the protan, not a common one (Section~\ref{sec:results:geometry}). Leave-one-color-out decoding shows \emph{where} the continuous-hue manifold breaks down: only hV4 supported above-chance interpolation in HC, replicating \citeA{brouwer2009}, and both CVD participants fell at or below chance there, most severely at the S-cone intermediate hues blue, purple, and magenta (Section~\ref{sec:results:loco}). This account aligns with prior characterizations of CVD as a multidimensional deformation of color space rather than a uniform attenuation \cite{boehm2014, emery2021}.

\subsection{A neurally grounded, individualizable correction filter}
We built each correction filter by inverting that participant's fitted two-component model to its stimulus-space pre-image (Methods~\S\ref{sec:methods:filter}) --- the displayed hue that, after the cortex distorts it, is perceived as the intended color. Across the eight hues this filter applied a mean rotation of $|\delta\theta| = 26.3^\circ$ for the deutan participant and $16.2^\circ$ for the protan participant.

The resulting filter is per-person and neurally grounded. To our knowledge, this is the first color-correction filter derived from an individual's cortical color representation rather than a retinal model. Building each filter from one participant's measured cortical color geometry rather than a fixed subtype-average spectral shift is the architectural advance over a population-average correction.

Fitting the behavioral and neural loss atoms independently revealed three benefits of grounding the correction in cortical geometry. First, the neural term recovered a confusion-axis direction the behavioral loss could not resolve: the protan behavioral-only fit did not beat the no-correction baseline ($\Delta L = +0.01$, 3 of 7 folds), whereas the RDM term resolved a clear rotation.
Second, it stabilized an already-informative solution: for the deutan participant the behavioral-only and combined fits shared the same argmin $(6^\circ, -42^\circ)$, and adding the neural term sharpened it without shifting it, lowering the boundary-saturation rate from $23\%$ to $9.3\%$.
Third, it tightened parameter stability in both participants, narrowing the resample interquartile range of the recovered parameters. The neural-only fit for the deutan participant independently returned the same direction ($\hat\beta_c = -26^\circ$), corroborating the behavioral estimate. A correction grounded in individual cortical geometry therefore differs in substance from a purely behavioral one.

The two fitted filters differ by subtype in the \emph{direction} of the dominant confusion-axis term ($\hat\beta_c = -42^\circ$ for the deutan participant versus $+24^\circ$ for the protan participant). Because $\theta_{\rm conf}$ is subtype-defined and the sample holds one deutan and one protan, this contrast is between-subtype. It rests on the recoverable features of the fit: across held-out reference sets, the two-component mechanism class and the sign of the dominant confusion-axis term are stable, whereas the per-axis magnitudes are not separately identifiable (0 of 6 recovery checks survive FDR correction). The contrast therefore rests on these stable directions, not on absolute parameter values.

A retina-based transform does not substitute for this per-person inverse, for a structural reason. The Machado class --- a fixed spectral shift along the single confusion axis, even when scaled by a cortical gain \cite{machado2009} --- stays confined to that one axis and therefore does not represent distortion in orthogonal directions. For the protan participant it even compresses three distinct displayed hues into a narrow arc, leaving no exact pre-image where that participant is most vulnerable (Appendix~A). Fitting the class confirmed the mismatch: its gain saturated near $g \approx 3$, implausibly implying overcompensation in observers who remain measurably impaired. But the structural limit alone rules it out as an exact correction (Section~\ref{sec:methods:candidates}).

\subsection{Filter evaluation}
The two evaluation domains diverged. Behaviorally, the individualized filter normalized the deutan participant's baseline discrimination deficit and, in the protan participant, normalized a trending confusion-axis pair (green--blue) while introducing no new deviation. The deployed accessibility filter, by contrast, produced two new deviant pairs in that participant. It did not, however, outperform the deployed filter in the deutan participant, where both removed the deficit. A behavioral advantage of the individualized filter over the deployed filter remains to be established with more participants. Neurally, the filter produced an inconsistent improvement. On the interpolation metric it increased adjacent accuracy in the deutan participant ($0.23 \to 0.31$) but decreased it in the protan participant ($0.14 \to 0.06$), and it reproduced a control-like geometry in neither. Filter efficacy therefore remains unproven and is the framework's immediate next test.

Neither filter reproduced the full control geometry, and the two participants moved in opposite directions. In the deutan participant the unfiltered baseline was the closest to HC of the three conditions, and both filters moved the geometry away from it. In the protan participant both filters moved the distorted baseline toward the controls, but comparably (deployed $\approx$ individualized), so the shift was not specific to the individualized filter. This early-visual geometry and the hV4 interpolation read-out lie at different levels of the visual hierarchy. Early areas encode the physical stimulus, so a recolored stimulus necessarily reshapes V1/V2 geometry, whereas hV4 reflects a later, perceptual stage. The two levels therefore need not change together, and their divergent filter effects are consistent with this hierarchical difference.

This evaluation rests on one deutan and one protan case. The primary-endpoint effect appeared in the deutan participant and reversed in the protan participant. Whether either pattern generalizes cannot be determined from two cases; replication in more individuals, within each subtype, is the primary requirement.

\subsection{Limitations}
Five considerations bound the proof-of-concept scope of these findings. First, the CVD sample is $N = 2$. Single-case statistics \cite{crawford1998} permit individual-level inference, but population-level claims about deutan or protan subtypes require replication with more participants. As established above, the sign of $\hat\beta_c$ is the recoverable quantity. Whether the between-subtype $\hat\beta_c$ difference ($-42^\circ$ vs.\ $+24^\circ$) is subtype-level or idiosyncratic awaits power-adequate replication. Establishing per-person rather than subtype-average correction specifically requires testing several individuals \emph{within} a single subtype --- the decisive follow-up experiment. Second, the stimulus set was confined to a single isoluminant, iso-chroma locus ($L^{*} = 75$, chroma $= 40$). The filter is therefore derived only along the hue circle at fixed lightness and saturation, and any future validation is similarly confined. It does not constrain how the correction generalizes to luminance- or saturation-varying stimuli or to naturalistic viewing --- gaps that denser stimulus sampling across these axes would close. Third, the HC reference cohort is small (seven controls), so the HC $\|\hat\beta\|$ leave-one-out distribution we report is a magnitude anchor for the CVD estimates, not a hypothesis test. A larger HC reference cohort would be required to arbitrate population-level distortion patterns. Fourth, we did not compute parameter-level bootstrap confidence intervals on $\hat\beta_s$ and $\hat\beta_c$, so the per-axis estimates are point values whose interval precision awaits the additional runs per participant that bootstrap resampling would require. Fifth, the fitting objective did not directly target the interpolation deficit. The interpolation (LOCO) signature was included among the candidate fitting terms, but for neither participant did it survive model selection; both winning combinations were the behavioral JND loss plus the neural $\Delta$RDM. The LOCO atom is defined in a forward-encoding channel space at hV4 --- a smaller voxel set on a different representational basis than the RDM --- which may have made it difficult to integrate stably with the geometry loss. A loss that stably incorporates the interpolation signature alongside the geometry is left to future work.

\subsection{Conclusion}
We first characterized the cortical color representation in two individuals with CVD: categorical discrimination of all displayed hues was preserved, whereas the continuous geometry relating them was selectively distorted, at a distinct cortical area in each participant. Building on this, we show as a proof-of-concept that each individual's cortical color geometry, measured with fMRI, can ground a stimulus-space correction obtained by inverting a fitted model of that individual's representational distortion. It illustrates, in two cases, that a color-vision deficiency need not present as a uniform retinal loss undone by a fixed subtype-average spectral shift. It can instead be an individually patterned, multidimensional distortion of the cortical color representation --- one structured enough to invert. Because the correction is the inverse of each individual's own fitted distortion, it is specific to that person rather than a fixed, population-average retinal transform.
However, in a follow-up two-person evaluation, the correction's effect was inconsistent across participants and metrics. Whether it can provide a generalizable improvement in color perception therefore remains to be established. Beyond this proof-of-concept, larger and systematic studies can quantify the cortical color distortions of CVD and provide a new class of individualized, cortically grounded filters.
